\chapter{Introducción}
\label{sec:introduccion}

El estudio de la exposición al ruido en los ambientes de trabajo constituye uno de los pilares fundamentales
de la higiene ocupacional. El ruido no controlado en las actividades laborales puede generar pérdidas auditivas
inducidas por ruido, alteraciones vasculares, estrés y aumento del riesgo de accidentes de trabajo.

El presente trabajo práctico tiene como objetivo llevar a cabo el relevamiento y la evaluación acústica de un
puesto de trabajo simulado de reparación de electrodomésticos, en el cual se utiliza una herramienta rotativa de
alta velocidad (minitorno Dremel). Para tal fin, se utiliza instrumental de medición digital integrador mediante
la aplicación Sound Analyzer v2.7 para Android, relevando el nivel sonoro instantáneo con ponderación A y
respuesta rápida ($L_{AF}$) y el nivel sonoro continuo equivalente ($L_{Aeq}$).

La operación ensayada presenta una dinámica de trabajo intermitente caracterizada por ciclos de $30\s$ de trabajo
activo y $5\s$ de reposicionamiento, evaluando dos regímenes de marcha diferenciados: $15.000\text{ RPM}$ para
operaciones de corte y $25.000\text{ RPM}$ para tareas de pulido y limpieza de óxido sobre componentes metálicos.

Asimismo, se evalúa la conformidad del puesto de trabajo frente a los límites permisibles establecidos por la
Ley N° 19.587 de Higiene y Seguridad en el Trabajo, el Decreto N° 351/79 y la Resolución N° 295/2003 del
Ministerio de Trabajo, Empleo y Seguridad Social (MTEySS). Finalmente, los datos relevados y los cálculos
de dosis proyectados a una jornada de 8 horas se vuelcan de forma completa en la plantilla oficial dictada
por la Resolución N° 85/2012 de la Superintendencia de Riesgos del Trabajo (SRT).
